% !TEX root = ../main.tex

\begin{abstract}[zh]
\addcontentsline{toc}{chapter}{摘\quad 要}
科研文献工作始于论文发现，并延续到理解、核验、保存和关联。论文修订会改变可用证据，研究兴趣也会随阅读发展。持续使用的文献助手需要明确论文版本、连接生成内容与原文，并在移动环境中以可检查的方式利用用户兴趣。

本文设计并实现 MNEME，一个面向持续文献理解的移动原生人工智能研究助手。十名学生和早期研究人员的访谈及后续设计分析确定了信息过载、论文找回、知识组织和来源可信度方面的需求。系统分离表示逻辑论文与观测版本，并将版本信息传递到检索片段、摘要和引用。章节感知检索与引用检查支持单篇论文摘要和自由问答；版本化行为事件与可解释排序形成可重放的兴趣更新。Android 客户端以本地缓存、可恢复事件和明确的内容来源状态连接论文简报、详情、问答和受限引文图。最终产品进一步支持连续问答、论文保存、公共分享、每周提醒和图谱筛选。

五名参与者的形成性测试中，推荐论文打开、来源核验和带引用问答均达到 5/5 完成。图谱探索观察形成了节点提示、持续选中、关系说明和“打开论文”等改进。

系统评价覆盖行为建模、证据受限问答和 Android 客户端。十一组受控行为轨迹中，模型的平均 nDCG@5 为 0.92，时效性基线为 0.92；平均倒数排名从 0.88 提高到 0.94。机制实验验证了暴露门控、负向反馈、单篇论文饱和和重复事件幂等性。十五个实时问答案例中，所有返回的答案均包含通过词汇来源检查的引用，十二个可回答案例均未被错误拒答。Android 资源评价在八个分别冷启动的模拟器会话中保留了 80 个启动与恢复计时和 320 个图谱计时；每次启动、恢复和图谱试验均达到预先定义的状态与内容条件。三个固定种子的首次初始化均显示了五篇实时论文，从提交到简报可见用时 28.46--41.73 秒，中位数为 36.50 秒。十组三事件同步批次在数据库重新打开和 HTTP~503 响应后保持待同步状态；隔离后端接收了 30 个唯一事件，并将随后提交的 30 个相同标识符识别为重复事件。系统还完成了从种子论文到五篇论文简报、论文详情、自由问答和多节点引文图的实时流程。最终集成流程进一步完成了论文保存与分享、两轮带来源问答、50 篇论文和 53 条有向引用的图谱筛选，以及由每周提醒进入对应的十篇论文简报；后端 716 项回归测试和 Android 32 项连接设备测试均全部通过。

本文提出并实现了一套连接论文版本、来源证据、行为记忆与移动交互的完整系统方法。实验结果表明，MNEME 能够支持基于来源的论文理解、可解释的兴趣更新和稳定的移动端研究流程，为持续文献管理与人工智能辅助阅读提供了可实现、可检验的系统方案。
\end{abstract}

{
\newfontface{\arial}{Arial}[BoldFont={Arial Bold},Scale=0.94]
\ctexset{chapter/format+={\arial}}
\begin{abstract}[en]
\addcontentsline{toc}{chapter}{ABSTRACT}
Scholarly literature work continues after paper discovery through understanding, verification, retention, and connection. Paper revisions change the available evidence, while research interests develop through reading. A mobile literature assistant must therefore preserve revision identity, source links, and inspectable interest evidence under changing network conditions.

This thesis designs and implements MNEME, a mobile-native AI research assistant for continuous literature understanding. Interviews with ten students and early-stage researchers, followed by design analysis, established requirements for information volume, paper recovery, knowledge organization, and source trust. The system separates logical papers from observed revisions and propagates revision identity to retrieved passages, summaries, and citations. Section-aware retrieval and citation checks support summaries and free-form questions. Versioned behavioral events and explainable ranking provide replayable interest updates. The Android client uses cached state, recoverable event delivery, and visible content origins across a five-paper briefing, paper detail, question answering, and a bounded citation graph. The final product also supports follow-up questions, paper retention, public sharing, weekly alerts, and graph filtering.

A formative study with five participants recorded 5/5 completion for opening a recommended paper, verifying a source, and using cited question answering. Graph observations led to clearer node guidance, persistent selection, relationship explanations, and a visible Open Paper action.

Across eleven controlled behavior traces, the model maintained mean nDCG@5 relative to a recency baseline (0.92 versus 0.92) while increasing mean reciprocal rank from 0.88 to 0.94. Mechanism tests confirmed exposure gating, negative feedback, per-paper saturation, and duplicate-event idempotency. In fifteen live question-answering cases, every returned answer contained citations that passed the lexical source check, and none of the twelve answerable cases was falsely refused. The Android resource evaluation retained 80 startup-and-recovery timings and 320 graph timings across eight separately cold-booted emulator sessions; every startup, recovery, and graph trial met its defined state and content checks. First initialization from each of three fixed seeds displayed five live papers in 28.46--41.73~s, with a median of 36.50~s. Ten three-event synchronization batches remained pending through database reopening and an HTTP~503 response; the isolated backend accepted all 30 unique events and identified all 30 subsequent replays as duplicates. The system also completed a live workflow from a seed paper to a five-paper briefing, paper detail, free-form question answering, and a multi-node citation graph. The final integrated flow additionally saved and shared the selected paper, completed two source-linked question turns, filtered a graph containing 50 papers and 53 directed citations, and opened the corresponding ten-paper weekly briefing from its alert; all 716 backend regression tests and all 32 Android connected-device tests passed.

This thesis presents and implements an integrated system that connects paper revision identity, source evidence, behavioral memory, and mobile interaction. The evaluation shows that MNEME supports source-grounded paper understanding, explainable interest adaptation, and a reliable mobile research workflow, providing an implementable and testable approach to continuous literature management and AI-assisted reading.
\end{abstract}
}
